\documentclass[11pt]{article}

\usepackage[T1]{fontenc}
\usepackage[margin=1.1in]{geometry}
\usepackage{amsmath,amssymb,amsthm}
\usepackage{enumitem}
\usepackage{booktabs}
\usepackage[colorlinks=true,linkcolor=blue,citecolor=blue,urlcolor=blue]{hyperref}

\newtheorem{theorem}{Theorem}[section]
\newtheorem{proposition}[theorem]{Proposition}
\newtheorem{lemma}[theorem]{Lemma}
\newtheorem{corollary}[theorem]{Corollary}
\newtheorem{definition}[theorem]{Definition}
\newtheorem{conjecture}[theorem]{Conjecture}
\newtheorem{question}[theorem]{Question}
\theoremstyle{remark}
\newtheorem{remark}[theorem]{Remark}
\newtheorem{example}[theorem]{Example}

\newcommand{\R}{\mathbb{R}}
\newcommand{\X}{X}
\newcommand{\Xhat}{\widehat{\X}}
\newcommand{\Mx}{M_{\X}}
\newcommand{\Mg}{M^{g}_{\X}}
\newcommand{\Rx}{\mathcal{R}_{\X}}
\newcommand{\Rg}{\mathcal{R}^{g}_{\X}}
\newcommand{\Hb}[1]{\mathcal{H}_{#1}}
\newcommand{\dist}{\operatorname{dist}}
\newcommand{\dX}{d_{\X}}
\newcommand{\conv}{\operatorname{conv}}
\newcommand{\hconv}{\operatorname{hconv}}
\newcommand{\inj}{\operatorname{inj}}
\newcommand{\sect}{\operatorname{sec}}
\newcommand{\sca}[2]{\left\langle #1,\, #2\right\rangle}
\newcommand{\E}{\mathbb{E}}
\newcommand{\HH}{\mathbb{H}}
\newcommand{\Sn}{\mathbb{S}^{n}}
\newcommand{\bdry}{\partial_{\infty}}

\title{Medial-axis reconstruction on complete Riemannian manifolds:\\
capture, drift, and the unconditional theorem\thanks{Final note of a
series generalising Theorem~6 of \cite{Bia}: \cite{R1} carried its
sufficiency half to Hadamard manifolds, \cite{R2} proved the full
characterisation in real hyperbolic space, \cite{R3} gave the faithful
curvature-free characterisation for every Hadamard manifold, \cite{Sph}
crossed into positive curvature, and \cite{Tri} organised the theory
into its two axes (compactness; the Busemann--Hessian sign). This note
removes the last analytic caveats and states the theorem on an
arbitrary complete Riemannian manifold. Drafted with the assistance of
Claude (Anthropic); all arguments were re-derived and checked by hand,
and the two mechanics of the capture theorem --- the rate-one identity
off the axis and the margin at first approach --- were verified
numerically on $\mathbb{S}^{2}$ against a non-circular convex contact
curve, whose focal region forces the foot to slide (60 random starts:
growth rate $1$ to five digits, margin deficit $<3\cdot10^{-6}$; file
\texttt{medial\_capture\_check.py}). Remaining errors are ours.}}
\author{}
\date{July 2026}

\begin{document}
\maketitle

\begin{abstract}
For a closed set $\X$ in a complete Riemannian manifold $M$, the
\emph{two-geodesic medial axis} $\Mg$ is the set of points admitting
two distinct minimising geodesics to $\X$ (the cut locus of $\X$), and
the reconstruction $\Rg=\bigcup_{a\in\Mg}B(a,\dX(a))$ is the union of
its open medial balls. Across the series, the generalisation of
Bia\l{}o\.zyt's Theorem~6 to Riemannian manifolds had been proved up to
two flagged gaps: the total-reconstruction theorem on compact manifolds
rested on an unproven extension of Lieutier's swallow property past cut
points \cite[Rem.~8.4]{Sph}, \cite[Rem.~6.5]{Tri}, and the Hadamard
characterisation for non-compact contact sets carried a radial-offset
caveat \cite[Thm.~3.4]{R3}. This note closes both, with no genericity,
analyticity, curvature, or spread hypothesis. The first input is a
\emph{capture theorem}: along any solution of the differential
inclusion $\Phi'\in-\conv U(\Phi)$, where $U$ is the (upper
semicontinuous, in general non-closed-graph-selectable) field of unit
directions of minimising geodesics to $\X$, the distance $\dX$ grows at
unit rate for exactly as long as the trajectory avoids $\Mg$; hence
every point of $M\setminus\X$ is either \emph{captured} --- it lies at
depth $\ge\dX(p)$ inside a medial ball of $\Mg$ --- or its trajectory
\emph{escapes} to infinity, which is impossible when $M$ is compact.
Consequently $\Rg=M\setminus\X$ for \emph{every} compact $M$ and every
closed $\X\neq\varnothing$ (even a single point). The subtlety that
defeated the margin arguments --- $\Mg$ is not closed, and the margin
is not monotone through multi-geodesic points --- is dissolved by
stopping at the \emph{first approach} to $\Mg$ rather than flowing
through it; focal and conjugate degenerations are ordinary points of
the inclusion. The second input is a horofunction formulation: calling
$\psi$ a \emph{medial limit function} if $\psi=\lim\,[d(\cdot,a_{k})
-\dX(a_{k})]$ locally uniformly with $a_{k}\in\Mg$, the swallow lemma
becomes three lines --- $\{\psi<0\}\subseteq\Rg$ --- with no anchoring
or pinching, and
\[
  \Rg=\bigcup\bigl\{\{\psi<0\}:\ \psi\ \text{a medial limit
  function}\bigr\}
\]
holds on \emph{every} complete $M$. On a Hadamard manifold ideal medial
limit functions are Busemann functions up to level, which resolves the
offset caveat of \cite{R3}: a drift of offset $c$ is exactly a medial
limit at level $c$, reconstructing $\{b_{\xi}<-c\}$ --- a statement,
no longer a gap. Together with the trichotomy
captured/drift-swallowed/escaping, this completes the programme: the
Euclidean, Hadamard, spherical, spread, and compact theorems are all
specialisations, and the remaining open problems are exactly the
faithful ($\Mg$-free) certificates on the escape side.
\end{abstract}

\tableofcontents

\section{Introduction}\label{sec:intro}

For a closed set $\X\neq\varnothing$ in a complete Riemannian manifold
$M$ write $\dX=\dist(\cdot,\X)$,
\[
  \Mx=\{a:\ a\ \text{has}\ \ge2\ \text{nearest points in}\ \X\},
  \qquad
  \Mg=\{a:\ a\ \text{has}\ \ge2\ \text{minimising geodesics to}\ \X\},
\]
for the two-\emph{point} and two-\emph{geodesic} medial axes
($\Mx\subseteq\Mg$; the endpoints of the two geodesics may coincide),
and
\[
  \Rx=\bigcup_{a\in\Mx}B\bigl(a,\dX(a)\bigr),\qquad
  \Rg=\bigcup_{a\in\Mg}B\bigl(a,\dX(a)\bigr)
\]
for the corresponding reconstructions by open medial balls.
Bia\l{}o\.zyt's Theorem~6 \cite{Bia} characterises $\Rx$ in $\R^{n}$
(where $\Mx=\Mg$); the companion notes carried the characterisation to
real hyperbolic space \cite{R2}, to every Hadamard manifold with a
faithful horoconvexity certificate \cite{R1,R3}, and into positive
curvature \cite{Sph}, and \cite{Tri} organised the results into two
independent axes: \emph{compactness} governs partial-versus-total,
the \emph{Busemann--Hessian sign} governs the detection flavour.

Two analytic gaps remained, both flagged explicitly in the series.

\begin{enumerate}[label=(G\arabic*)]
\item\label{G1} \emph{The containment step on compact manifolds.} The
total-reconstruction theorem $\Rg=M\setminus\X$ for compact $M$
\cite[Thm.~8.3]{Sph}, \cite[Thm.~6.4]{Tri} used the ascending
generalised-gradient flow of $\dX$ and needed the starting point to
remain inside the moving free ball. The elementary constant-margin
identity proves this only while the reach stays below the cut distance
of the current foot; past it, the argument invoked ``Lieutier's swallow
property extending past cut points --- expected, and unconditional for
generic or analytic metrics, but not delivered by the margin identity
alone'' \cite[Rem.~8.4]{Sph}, \cite[Rem.~6.5]{Tri}.
\item\label{G2} \emph{The radial offset for non-compact contact sets.}
In the Hadamard characterisation for arbitrary closed $\X$
\cite[Thm.~3.4]{R3}, a purely asymptotic drift in an inhomogeneous
manifold could carry a nonvanishing radial offset, and then only the
deep part of the defective horoball was proved reconstructed.
\end{enumerate}

This note closes \ref{G1} and \ref{G2} and states the resulting
theorem in full generality: \emph{every} complete Riemannian manifold,
\emph{every} closed set. No curvature bound, genericity, analyticity,
spread, or cardinality hypothesis is used anywhere.

\subsection*{The two inputs}

\emph{Capture (\S\ref{sec:capture}).} The obstruction in \ref{G1} was
never the existence of the ascending flow; it was the bookkeeping of
the margin through multi-geodesic points, where the flow turns and the
margin genuinely need not be monotone (Remark~\ref{rem:notmonotone}).
The repair is to reverse the roles: follow \emph{any} trajectory of
the differential inclusion driven by the minimising-geodesic
directions, and \emph{stop at the first approach to $\Mg$}. Off $\Mg$
the direction field is single-valued, the trajectory is a unit-speed
curve along which $\dX$ grows at rate exactly one --- through focal and
conjugate points without comment --- so
\[
  \dX(\Phi_{t})=\dX(p)+t,\qquad d(p,\Phi_{t})\le t,
\]
and the starting point sits at depth $\ge\dX(p)$ in every moving ball.
The trajectory therefore either approaches $\Mg$ in finite time ---
\emph{capture}: some $a\in\Mg$ has $d(p,a)<\dX(a)$, i.e.\ $p\in\Rg$ ---
or avoids $\Mg$ forever and $\dX\to\infty$ along it --- \emph{escape},
impossible on a compact manifold. The fact that $\Mg$ is not closed
(minimising geodesics can merge in the limit: focal degenerations) is
harmless, because approach suffices: the margin is continuous. This is
Theorem~\ref{thm:capture}, and it makes the compact theorem
unconditional (Corollary~\ref{cor:compact}), for every closed
$\X\neq\varnothing$ --- including $|\X|=1$, which the two-point axis
cannot see at all.

\emph{Drift (\S\ref{sec:drift}).} Capture cannot be the whole story on
a non-compact manifold: a point can be reconstructed by medial balls
whose centres run to infinity while its own trajectory escapes (the
open half-planes of two points in $\R^{2}$ are the primal example).
The right limit object costs nothing in a complete manifold: $M$ is
proper (Hopf--Rinow), so the functions
$\psi_{a}=d(\cdot,a)-\dX(a)$ --- the \emph{free-ball functions} of
medial centres, normalised by their own reach, with no basepoint ---
are precompact in the topology of locally uniform convergence along any
sequence that reconstructs a fixed point. Calling any such limit a
\emph{medial limit function}, the swallow lemma of \cite{R3} becomes
three curvature-free lines with no anchoring hypothesis:
$\{\psi<0\}\subseteq\Rg$ for every medial limit function $\psi$
(Lemma~\ref{lem:swallow}). The characterisation on an arbitrary
complete manifold follows (Theorem~\ref{thm:char}), together with the
trichotomy: every $p\in M\setminus\X$ is captured, drift-swallowed, or
escaping-unswallowed, and $\Rg$ is exactly the first two classes
(Theorem~\ref{thm:trichotomy}). On a Hadamard manifold the ideal medial
limit functions are Busemann functions up to an additive level, and the
offset caveat \ref{G2} dissolves into a definition: a drift with
asymptotic offset $c$ \emph{is} a medial limit function at level $c$
and reconstructs $\{b_{\xi}<-c\}$ in full
(Remark~\ref{rem:offset}).

\subsection*{What the completed theory says}

Section~\ref{sec:spec} re-derives every earlier result as a
specialisation and updates the $2\times2$ of \cite{Tri}: the compact
row is unconditionally total for $\Rg$; the non-compact row is governed
by the medial limit functions, whose faithful ($\Mg$-free)
certificates are known exactly on the geometries where the series
established them --- $\conv$ in $\R^{n}$ \cite{Bia}, $\hconv$ on
Hadamard manifolds \cite{R3}, vacuous on compact $M$ --- and open in
mixed curvature. Those certificates, and the quantitative
cut-shielding threshold for the two-point axis \cite[Rem.~6.6]{Tri},
are the honest residue of the programme
(\S\ref{sec:open}).

\section{Preliminaries}\label{sec:prelim}

Throughout, $M$ is a complete Riemannian manifold (hence a proper
geodesic metric space, by Hopf--Rinow), $\X\subseteq M$ is closed with
$\varnothing\neq\X\neq M$, and $\dX=\dist(\cdot,\X)$. All geodesics are
unit-speed. For $x\in M\setminus\X$ let
\[
  U(x)=\bigl\{\gamma'(0):\ \gamma:[0,\dX(x)]\to M\ \text{minimising
  geodesic},\ \gamma(0)=x,\ \gamma(\dX(x))\in\X\bigr\}
  \subseteq T^{1}_{x}M ,
\]
the set of unit initial directions of minimising geodesics to $\X$.
Since a geodesic is determined by its initial vector, distinct
minimising geodesics have distinct initial directions, so
\[
  \Mg=\{x\in M\setminus\X:\ |U(x)|\ge2\}.
\]
We write $\Sigma:=\Mg$ when we use it as a stopping set, and
$r(a)=\dX(a)$ for the reach of a centre.

\begin{lemma}[Direction field: compactness and semicontinuity]
\label{lem:usc}
For every $x\in M\setminus\X$ the set $U(x)$ is nonempty and compact,
and the multimap $x\mapsto U(x)$ is upper semicontinuous: if
$x_{k}\to x$ in $M\setminus\X$ and $u_{k}\in U(x_{k})$, $u_{k}\to u$
(in $TM$), then $u\in U(x)$. In particular $U$ is single-valued and
continuous at every point of $(M\setminus\X)\setminus\Sigma$.
\end{lemma}

\begin{proof}
Nonemptiness and compactness: minimisers exist by properness, and a
limit of initial vectors of minimising geodesics from the same $x$ is
the initial vector of a limit geodesic, which minimises by continuity
of $\dX$ and of lengths. Upper semicontinuity: let
$\gamma_{k}$ be the minimising geodesics with
$\gamma_{k}'(0)=u_{k}$. By smooth dependence of geodesics on initial
conditions, $\gamma_{k}\to\gamma$ locally uniformly where
$\gamma'(0)=u$; the lengths $\dX(x_{k})\to\dX(x)$, and the endpoints
$\gamma_{k}(\dX(x_{k}))\in\X$ converge to $\gamma(\dX(x))\in\X$ ($\X$
closed). So $\gamma$ is a geodesic from $x$ to $\X$ of length
$\dX(x)$: minimising, i.e.\ $u\in U(x)$. At a point with
$U(x)=\{u\}$, upper semicontinuity of a compact-valued multimap into a
singleton forces continuity of every selection near $x$ --- note this
is continuity \emph{at} $x$ only; $U$ may be multivalued arbitrarily
close by.
\end{proof}

\begin{lemma}[First variation of $\dX$; {\cite[\S3]{MM}}]
\label{lem:firstvar}
For $x\in M\setminus\X$ and $w\in T_{x}M$, the one-sided directional
derivative of $\dX$ exists and equals
\[
  \frac{d}{ds}\Big|_{s=0^{+}}\dX\bigl(\exp_{x}(sw)\bigr)
  =\min_{u\in U(x)}\sca{w}{-u}.
\]
\end{lemma}

\begin{proof}[Proof sketch]
``$\le$'': fix $u\in U(x)$ with foot $y=\gamma(\dX(x))$ and vary
$\gamma$ through curves from $\exp_{x}(sw)$ to $y$; the first variation
of arc length gives
$\dX(\exp_{x}(sw))\le d(\exp_{x}(sw),y)\le\dX(x)+s\sca{w}{-u}+O(s^{2})$.
``$\ge$'': for $s_{k}\downarrow0$ pick $u_{k}\in U(\exp_{x}(s_{k}w))$;
by Lemma~\ref{lem:usc} a subsequence converges to some $u\in U(x)$, and
the reverse comparison along the geodesics with directions $u_{k}$
gives the matching lower bound with this $u$, hence with the minimum.
Details (with the semiconcavity of $\dX$ on $M\setminus\X$, which we do
not need) are in Mantegazza--Mennucci \cite{MM}.
\end{proof}

\paragraph{The descent inclusion.} On the open set $M\setminus\X$
define the multimap
\[
  G(x)\ :=\ \conv\bigl(-U(x)\bigr)\ \subseteq\ T_{x}M .
\]
By Lemma~\ref{lem:usc}, $G$ is upper semicontinuous with nonempty
compact convex values contained in the closed unit ball. A
\emph{capture curve} from $p\in M\setminus\X$ is a locally Lipschitz
curve $\Phi:[0,T_{\max})\to M\setminus\X$ with $\Phi_{0}=p$ and
\begin{equation}\label{eq:incl}
  \Phi'(t)\in G(\Phi_{t})\qquad\text{for a.e.\ }t ,
\end{equation}
maximally extended in $M\setminus\X$. Local existence of solutions of
\eqref{eq:incl}, and extension as long as the trajectory stays in a
compact subset of the open domain, are the standard theory of
differential inclusions with upper semicontinuous convex compact
right-hand side (Aubin--Cellina \cite[\S2.1]{AC}), applied in charts
and concatenated; $|\Phi'|\le1$ a.e., so trajectories are
$1$-Lipschitz. We emphasise that \emph{no} uniqueness, canonicity, or
semiconcave-gradient selection is needed: any solution serves. (At a
point of $\Sigma$ the inclusion may branch; off $\Sigma$ it is the
classical ODE of the field $-u$.) When $U$ is single-valued along a
trajectory the curve simply extends the current minimising geodesic
beyond $x$, switching fields continuously as the foot slides.

\paragraph{Free-ball functions.} For $a\in M\setminus\X$ set
\[
  \psi_{a}:=d(\cdot,a)-\dX(a)\ :\ M\to\R ,
\]
a $1$-Lipschitz function, negative exactly on the open free ball
$B(a,\dX(a))$, with $\psi_{a}\ge0$ on $\X$. It is normalised by the
reach --- $\min\psi_{a}=-\dX(a)$ at $a$ --- rather than by a basepoint;
this intrinsic normalisation is what makes the limits below
basepoint-free.

\begin{definition}[Medial limit functions]\label{def:mlf}
A function $\psi:M\to\R$ is a \emph{medial limit function} of $\X$ if
there are $a_{k}\in\Mg$ with $\psi_{a_{k}}\to\psi$ locally uniformly.
It is \emph{finite} if $(a_{k})$ has a bounded subsequence and
\emph{ideal} otherwise (then $\psi$ is a horofunction-type limit:
$d(a_{k},\cdot)\to\infty$ locally uniformly while
$\psi$ stays finite). Every $\psi$ is $1$-Lipschitz, and
$\psi\ge0$ on $\X$ (each $\psi_{a_{k}}$ is), so $\{\psi<0\}$ is an open
set disjoint from $\X$: the \emph{free horoball} of $\psi$. For finite
$\psi$ along $a_{k}\to a$: if $a\in\Mg$ then $\psi=\psi_{a}$ is a
free-ball function; in general ($\Mg$ is not closed) $\psi=\psi_{a}$
with $a$ a focal-degenerate limit of medial centres --- still usable,
by Lemma~\ref{lem:swallow}, which never needs the limit centre itself.
\end{definition}

\begin{remark}[Relation to the horofunction compactification]
\label{rem:horo}
$M$ proper makes $C(M)$ with the topology of locally uniform
convergence a countably compact home for $1$-Lipschitz functions
normalised at a point (Arzel\`a--Ascoli). The classical horofunction
compactification \cite{Gro,BH} normalises $d(\cdot,a_{k})$ by a
basepoint value; our $\psi$'s differ from it by the scalar
$d(o,a_{k})-\dX(a_{k})$, which is exactly the \emph{offset} bookkeeping
of \cite[Thm.~3.4]{R3}. Working with reach-normalised limits absorbs
the offset into the function itself; see Remark~\ref{rem:offset}. On a
Hadamard manifold every ideal medial limit function is a Busemann
function plus a constant (horofunctions of a proper $\mathrm{CAT}(0)$
space are Busemann \cite[II.8.13]{BH}), so ideal medial limit
functions are exactly ``Busemann function $b_{\xi}$ $+$ level $c$''.
\end{remark}

\section{The capture theorem}\label{sec:capture}

\begin{theorem}[Capture or escape]\label{thm:capture}
Let $M$ be complete, $\X$ closed, $p\in M\setminus\X$, and let $\Phi$
be any capture curve from $p$. Let
\[
  T\ :=\ \sup\bigl\{t\in[0,T_{\max}):\ \Phi_{s}\notin\Sigma\ \ \forall
  s\in[0,t]\bigr\}.
\]
Then on $[0,T)$,
\begin{equation}\label{eq:rateone}
  \dX(\Phi_{t})=\dX(p)+t,
  \qquad
  d(p,\Phi_{t})\le t,
  \qquad\text{hence}\qquad
  \dX(\Phi_{t})-d(p,\Phi_{t})\ \ge\ \dX(p),
\end{equation}
and exactly one of the following holds.
\begin{enumerate}[label=\textup{(\roman*)}]
\item\emph{(Capture.)} $T<T_{\max}$. Then there exist $a\in\Mg$
arbitrarily close to $\Phi_{T}$ \textup{(}$a=\Phi_{T}$ itself if
$\Phi_{T}\in\Sigma$\textup{)} with
\[
  \dX(a)-d(p,a)\ \ge\ \dX(p)-\varepsilon
  \qquad(\varepsilon>0\ \text{arbitrary}),
\]
so $p\in B(a,\dX(a))$ for such $a$: \quad $p\in\Rg$, at depth
arbitrarily close to $\dX(p)$, and $T\le\operatorname{diam}(M)-\dX(p)$.
\item\emph{(Escape.)} $T=T_{\max}$. Then $T_{\max}=\infty$, the
trajectory never meets $\Sigma$, $\dX(\Phi_{t})=\dX(p)+t\to\infty$, and
$\Phi_{t}$ leaves every compact subset of $M$. This is possible only
if $M$ is non-compact.
\end{enumerate}
\end{theorem}

\begin{proof}
\emph{The rate-one identity.} Fix $t\in[0,T)$; on $[0,t]$ the
trajectory avoids $\Sigma$, so $U(\Phi_{s})=\{u(\Phi_{s})\}$ is a
singleton for every $s\le t$ and \eqref{eq:incl} reads
$\Phi'(s)=-u(\Phi_{s})$ at a.e.\ $s$. The map $h(s):=\dX(\Phi_{s})$ is
Lipschitz. Let $s$ be any point where both $\Phi'(s)$ exists (with the
inclusion valid) and $h'(s)$ exists --- almost every $s\in[0,t]$. Write
$v=\Phi'(s)=-u(\Phi_{s})$, a unit vector. In a chart around $\Phi_{s}$,
differentiability gives
$d\bigl(\Phi_{s+\sigma},\exp_{\Phi_{s}}(\sigma v)\bigr)=o(\sigma)$ as
$\sigma\downarrow0$, and $\dX$ is $1$-Lipschitz, so
\[
  h(s+\sigma)=\dX\bigl(\exp_{\Phi_{s}}(\sigma v)\bigr)+o(\sigma)
  =h(s)+\sigma\min_{u\in U(\Phi_{s})}\sca{v}{-u}+o(\sigma)
  =h(s)+\sigma\,\sca{-u}{-u}+o(\sigma),
\]
by Lemma~\ref{lem:firstvar} and single-valuedness. Hence
$h'(s)=|u|^{2}=1$ at a.e.\ $s\in[0,t]$, and integrating,
$h(t)=\dX(p)+t$. The bound $d(p,\Phi_{t})\le t$ is the $1$-Lipschitz
property of the trajectory, and the margin inequality in
\eqref{eq:rateone} is their difference. Note that no closedness of
$\Sigma$, and no structure at focal or conjugate points, was used: a
point of $(M\setminus\X)\setminus\Sigma$ in the closure of $\Sigma$ is
simply a point where $U$ is single-valued, and the identity passes
through it.

\emph{Dichotomy.} Suppose first $T<T_{\max}$. If $\Phi_{T}\in\Sigma$,
continuity of $\dX(\Phi_{\cdot})-d(p,\Phi_{\cdot})$ and
\eqref{eq:rateone} on $[0,T)$ give
$\dX(\Phi_{T})-d(p,\Phi_{T})\ge\dX(p)>0$, so $a=\Phi_{T}\in\Mg$
contains $p$ in its open medial ball at depth $\ge\dX(p)$. If
$\Phi_{T}\notin\Sigma$, then by definition of $T$ as a supremum, every
interval $(T,T+\delta]\subseteq[0,T_{\max})$ contains a time
$s_{\delta}$ with $\Phi_{s_{\delta}}\in\Sigma$; letting
$\delta\downarrow0$ and using continuity of the margin,
$a=\Phi_{s_{\delta}}\in\Mg$ satisfies
$\dX(a)-d(p,a)\ge\dX(p)-\varepsilon$ for $\delta$ small. In both cases
$p\in B(a,\dX(a))$ with $a\in\Mg$, i.e.\ $p\in\Rg$; and
$\dX\le\operatorname{diam}(M)$ bounds $T$ via \eqref{eq:rateone}.

Now suppose $T=T_{\max}$: the trajectory never meets $\Sigma$, and
\eqref{eq:rateone} holds on all of $[0,T_{\max})$. If
$T_{\max}<\infty$, then $\Phi([0,T_{\max}))\subseteq
\bar B(p,T_{\max})\cap\{\dX\ge\dX(p)\}$, a compact subset of the open
set $M\setminus\X$ (properness of $M$); by the extension property of
solutions this contradicts maximality. So $T_{\max}=\infty$ and
$\dX(\Phi_{t})=\dX(p)+t\to\infty$; since $\dX$ is bounded on compact
sets, $\Phi_{t}$ leaves every compact set. On compact $M$, $\dX\le
\operatorname{diam}(M)$ makes this impossible, so case (i) occurs.
\end{proof}

\begin{corollary}[Total reconstruction on compact manifolds,
unconditional]\label{cor:compact}
Let $M$ be a compact Riemannian manifold and $\X\subseteq M$ closed,
$\varnothing\neq\X\neq M$. Then
\[
  \boxed{\ \Rg=M\setminus\X\ },
\]
and quantitatively: every $p\in M\setminus\X$ lies at depth
$\ge\dX(p)-\varepsilon$ (any $\varepsilon>0$) inside a medial ball of
$\Mg$. No curvature, genericity, analyticity, spread, or cardinality
hypothesis is required; in particular $|\X|=1$ is allowed, where the
two-point axis has $\Mx=\varnothing$ but the two-geodesic axis still
reconstructs everything. This closes the containment caveat of
\cite[Rem.~8.4]{Sph} and \cite[Rem.~6.5]{Tri}: Theorem~8.3 of
\cite{Sph} and Theorem~6.4 of \cite{Tri} hold as stated, with
``resting on Lieutier's property'' deleted.
\end{corollary}

\begin{proof}
$\Rg\subseteq M\setminus\X$ since medial balls are free. Conversely,
if $p\in\Mg$ then $p\in B(p,\dX(p))\subseteq\Rg$; otherwise start a
capture curve at $p$ (local existence, \S\ref{sec:prelim}); by
Theorem~\ref{thm:capture} escape is impossible on compact $M$, so $p$
is captured.
\end{proof}

\begin{corollary}[Two-point axis under spread; {cf.\
\cite[Thm.~6.1]{Sph}}]\label{cor:spread}
If $M$ is compact and $\dX<\inj(M)$ on $M\setminus\X$ --- e.g.\ if
$\X$ is $\delta$-dense with $\delta<\inj(M)$ --- then
$\Rx=M\setminus\X$ for the two-point axis as well.
\end{corollary}

\begin{proof}
By Corollary~\ref{cor:compact} every $p$ lies in $B(a,\dX(a))$ for
some $a\in\Mg$. Two distinct minimising geodesics from $a$ with the
same endpoint $x_{0}$ would put $x_{0}$ in the cut locus of $a$, giving
$\dX(a)=d(a,x_{0})\ge\inj(M)$, contradiction. So the two geodesics of
$a$ have distinct endpoints: $a\in\Mx$.
\end{proof}

\begin{remark}[Why stopping at $\Sigma$, not flowing through it, is the
correct formulation]\label{rem:notmonotone}
Two natural strengthenings of Theorem~\ref{thm:capture} are false, and
identifying their failure is what shaped the proof. First, \emph{the
margin $\dX(\Phi)-d(p,\Phi)$ is not monotone through multi-geodesic
arcs}, already in $\R^{2}$: for $\X=\{(0,0),(0,2)\}$ and $p=(1,0.9)$,
the Lieutier flow reaches the bisector $\{y=1\}$ and slides along it,
and the margin decreases strictly (from $\approx1.345$ toward the
limit $1$, the depth of $p$ in the limiting half-plane). Second, even
the quadratic margin $\dX^{2}-d(p,\cdot)^{2}$, which \emph{is}
monotone along unique-foot and bisector arcs in $\R^{n}$, fails at
turning points where the active contact set changes (a three-point
configuration in $\R^{2}$ suffices: $\X=\{(0,0),(0,2),(10,5)\}$, with
the flow from $p=(1,0.9)$ turning at $\approx(5.75,1)$, where the
derivative of the quadratic margin is negative). Both phenomena live
\emph{on} $\Sigma$; on its complement the rate-one identity
\eqref{eq:rateone} is exact. Since reaching $\Sigma$ is the goal, not
an obstacle, the correct move is to stop there --- capture is
achieved, with the margin intact, \emph{before} the first
multi-geodesic time. The failure of closedness of $\Sigma$ (focal
degenerations, e.g.\ the evolute of an ellipse in $\R^{2}$, or the
conjugate-cusp endpoints of an ellipsoid's cut segment) is equally
harmless: such points are unique-geodesic points, the field $-u$ is
continuous at them (Lemma~\ref{lem:usc}), and the trajectory passes
through with $\dX$ still growing at rate one, the foot sliding
continuously.
\end{remark}

\begin{remark}[Relation to Lieutier's flow]\label{rem:lieutier}
For semiconcave $\dX$ one can canonically select the minimal-norm
element of $G$ (Lieutier's generalised gradient \cite{Lie}, in
$\R^{n}$; \cite{ACa,MM} for the semiconcavity). Our argument uses the
opposite regime: it only ever analyses the inclusion \emph{off}
$\Sigma$, where every selection agrees with the naive one (extend the
unique minimising geodesic), and it terminates at $\Sigma$, where
Lieutier's flow is designed to continue. This is why no analogue of
the swallow property past cut points is needed: the trajectory is
stopped at first approach, and the reconstruction is read off the
margin accumulated strictly before it. The price --- capture produces
a centre only \emph{near} the first $\Sigma$-approach when
$\Phi_{T}\notin\Sigma$ --- costs an arbitrarily small $\varepsilon$ of
depth, invisible to the open-ball statement.
\end{remark}

\section{Drift: the swallow, the characterisation, the trichotomy}
\label{sec:drift}

Capture is complete on compact manifolds, but on a non-compact $M$ a
point can be reconstructed by medial balls whose centres go to
infinity while its own capture curve escapes. The primal example is
$\X=\{x_{0},x_{1}\}\subseteq\R^{2}$: a point $p$ in an open half-plane
bounded by the line $x_{0}x_{1}$, placed beyond $x_{0}$ at an obtuse
angle, has a capture curve that runs straight to infinity with the
single foot $x_{0}$ --- yet $p$ is swallowed by the bisector balls
$B(a_{k},r(a_{k}))$, $a_{k}\to\infty$. The limit of those balls is a
half-plane, and the correct general home for such limits is the space
of medial limit functions of Definition~\ref{def:mlf}.

\begin{lemma}[Swallow, general and offset-free]\label{lem:swallow}
Let $M$ be any complete Riemannian manifold, $\X$ closed, and $\psi$ a
medial limit function of $\X$. Then
\[
  \{\psi<0\}\ \subseteq\ \Rg .
\]
If the witnessing centres lie in $\Mx$, then $\{\psi<0\}\subseteq\Rx$.
\end{lemma}

\begin{proof}
Let $\psi=\lim\psi_{a_{k}}$ locally uniformly, $a_{k}\in\Mg$, and let
$q\in\{\psi<0\}$. Then $\psi_{a_{k}}(q)\to\psi(q)<0$, so for large $k$
\[
  d(q,a_{k})-\dX(a_{k})=\psi_{a_{k}}(q)<0,
\]
i.e.\ $q\in B(a_{k},\dX(a_{k}))$ with $a_{k}\in\Mg$: $q\in\Rg$.
\end{proof}

No anchoring of feet, no curvature bound, no pinching, no offset
hypothesis: the entire content of the swallow lemmas of
\cite[Lem.~2.2]{R3} (and its constant-margin spherical form, \cite[Rem.~4.2]{Sph}) is absorbed into the
normalisation of $\psi_{a}$ by the reach.

\begin{lemma}[Compactness of reconstructing sequences]\label{lem:cpt}
Let $p\in\Rg$, and let $a_{k}\in\Mg$ be any sequence with
$\limsup_{k}\bigl[\dX(a_{k})-d(p,a_{k})\bigr]>0$ (e.g.\ a sequence of
centres whose balls contain $p$ with margin bounded away from $0$; a
single ball repeated qualifies). Then a subsequence of
$(\psi_{a_{k}})$ converges locally uniformly to a medial limit
function $\psi$ with $\psi(p)<0$.
\end{lemma}

\begin{proof}
The functions $\psi_{a_{k}}$ are $1$-Lipschitz and
$\psi_{a_{k}}(p)=d(p,a_{k})-\dX(a_{k})$ is bounded: below by
$-\dX(p)$ (since $\dX(a_{k})\le\dX(p)+d(p,a_{k})$) and above by $0$
along the qualifying subsequence. By Arzel\`a--Ascoli on the proper
space $M$, a further subsequence converges locally uniformly; the
limit is a medial limit function, and
$\psi(p)=\lim\psi_{a_{k}}(p)\le-\limsup[\dX(a_{k})-d(p,a_{k})]<0$.
\end{proof}

\begin{theorem}[Characterisation on an arbitrary complete manifold]
\label{thm:char}
For every complete Riemannian manifold $M$ and every closed
$\X\subseteq M$, $\varnothing\neq\X\neq M$,
\[
  \boxed{\ \Rg\ =\ \bigcup\bigl\{\{\psi<0\}:\ \psi\ \text{a medial
  limit function of}\ \X\bigr\}\ }.
\]
The same identity holds for $\Rx$ and $\Mx$ in place of $\Rg$ and
$\Mg$. (Every finite medial limit function is itself a free-ball
function $\psi_{a}$, with $a$ in the closure of $\Mg$ in
$M\setminus\X$; the sublevel set of such a limit is still swallowed,
by the lemma, even when $a\notin\Mg$.)
\end{theorem}

\begin{proof}
``$\supseteq$'' is Lemma~\ref{lem:swallow}. ``$\subseteq$'': if
$p\in\Rg$ then $p\in B(a,\dX(a))$ for some $a\in\Mg$, and
$\psi_{a}$ --- a medial limit function via the constant sequence ---
has $\psi_{a}(p)<0$.
\end{proof}

The identity is deliberately economical: its force is not the set
equality (the finite members alone give a tautology) but the pair of
facts it packages. \emph{First}, the ideal members reconstruct their
\emph{entire} open sublevel set --- the statement that in
\cite{R2,R3} required anchoring, pinching, or bisector witnesses, and
now holds bare. \emph{Second}, combined with
Theorem~\ref{thm:capture} it yields the trichotomy that generalises
the hull/drift/ray-escape structure of \cite[Thm.~2.1]{R3} to
arbitrary complete manifolds.

\begin{theorem}[Trichotomy on a complete manifold]\label{thm:trichotomy}
Let $M$ be complete, $\X$ closed, $p\in M\setminus\X$. Exactly one of
the following holds.
\begin{enumerate}[label=\textup{(\arabic*)}]
\item\emph{(Captured.)} Some capture curve from $p$ approaches $\Mg$
in finite time. Then $p\in\Rg$, at depth arbitrarily close to
$\dX(p)$.
\item\emph{(Swallowed, uncaptured.)} No capture curve from $p$
approaches $\Mg$ in finite time, but $\psi(p)<0$ for some medial limit
function $\psi$. Then $p\in\Rg$: reconstructed, but only by centres
$p$'s own ascent cannot find --- a remote ball, or a drifting sequence.
\item\emph{(Escaping.)} Neither. Then $p\notin\Rg$, and every capture
curve from $p$ escapes to infinity with $\dX(\Phi_{t})=\dX(p)+t$.
\end{enumerate}
In particular $\Rg=\{p:\ \textup{(1) or (2)}\}$, and on compact $M$
class \textup{(1)} is all of $M\setminus\X$.
\end{theorem}

\begin{proof}
The classes are exclusive and exhaustive by construction. (1)
$\Rightarrow p\in\Rg$ is Theorem~\ref{thm:capture}(i); (2)
$\Rightarrow p\in\Rg$ is Lemma~\ref{lem:swallow}. Conversely, if
$p\in\Rg$ then $p\in B(a,\dX(a))$ for some $a\in\Mg$ and the finite
medial limit function $\psi_{a}$ has $\psi_{a}(p)<0$, so $p$ lies in
class (2) whenever it is not in class (1) --- never in class (3).
Finally, a point not in class (1) has every capture curve uncaptured,
hence escaping, by the dichotomy of Theorem~\ref{thm:capture}.
\end{proof}

\begin{remark}[The witnesses of class (2)]\label{rem:class2}
Both kinds of witness genuinely occur, already for
$\X=\{x_{0},x_{1}\}\subseteq\R^{2}$ and $p$ beyond $x_{0}$ in an open
half-plane, at an obtuse angle to the bisector: every capture curve
from $p$ runs straight to infinity with the single foot $x_{0}$, yet
$p$ lies in a single \emph{finite} bisector ball (and also in the
ideal half-plane limit of the bisector balls). Class (2) is thus
``reconstruction is non-local for $p$'': the certificate cannot be
found by ascending from $p$, which is exactly the phenomenon that
separates the non-compact theory from the compact one and forces the
faithful certificates of \cite{Bia,R2,R3} to quantify over supporting
half-spaces or horoballs rather than over flow lines.
\end{remark}

\begin{remark}[The offset caveat of \cite{R3}, resolved]
\label{rem:offset}
Let $M$ be Hadamard. Ideal medial limit functions are exactly
$b_{\xi}+c$ with $\xi\in\bdry M$ and $c\in\R$
(Remark~\ref{rem:horo}), normalised so that $b_{\xi}$ vanishes on the
horosphere of the maximal free horoball $\Hb{}$ at $\xi$; necessarily
$c\ge0$ when $\X$ meets every neighbourhood of that horosphere. Call
$\Hb{}$ \emph{defective at level $c$} if $b_{\xi}+c$ is a medial limit
function; then Lemma~\ref{lem:swallow} reconstructs
$\{b_{\xi}<-c\}$ --- the horoball shrunk by $c$. The
``radial offset'' $\mathrm{off}_{k}=d(o,a_{k})+b_{\xi}(a_{k})$ of
\cite[Thm.~3.4]{R3} is precisely the discrepancy
$\psi_{a_{k}}-(b_{\xi}\circ\,\cdot\,)$ evaluated along the drift:
a drifting sequence with $\mathrm{off}_{k}\to c$ is a medial limit at
level $c$. The single unresolved case of \cite[Thm.~3.4]{R3} --- a
purely asymptotic drift with non-vanishing offset in an inhomogeneous
manifold, where only the deep part
$\{\tau_{p}>\limsup\mathrm{off}_{k}\}$ was reconstructed --- is
therefore not a gap in the theorem but the exact content of the level:
the drift \emph{is} a level-$c$ certificate and reconstructs exactly
$\{b_{\xi}<-c\}$, while the full-horoball statement is the assertion
that a level-$0$ certificate exists. Under any of the conditions of
\cite[Thm.~3.4]{R3} (pinching, flat case, a bisector witness on the
horosphere --- hence in every model of the series) level $0$ is
achieved and the horoball is reconstructed in full; the
$\hconv$-certificate of \cite[Thm.~3.3--3.4]{R3} detects
defectiveness at \emph{some} level unconditionally. The
characterisation of the exact level
\[
  c(\xi)\ :=\ \inf\{c\ge0:\ b_{\xi}+c\ \text{is a medial limit
  function}\}
\]
from $\X$ alone, in an inhomogeneous manifold, is the sharpened form
of the remaining question, and it is a question about $\X$ and the
horosphere geometry, not about the medial axis.
\end{remark}

\section{Specialisations: the series re-derived}\label{sec:spec}

\paragraph{Euclidean space \cite{Bia}.} In $\R^{n}$ geodesics between
two points are unique, so $\Mg=\Mx$ and $\Rg=\Rx$. Ideal medial limit
functions are affine: $|y-a_{k}|-r_{k}$ with $a_{k}\to\infty$
subconverges to $y\mapsto-\sca{y}{u}-c$, whose sublevel set is an open
half-space. Theorem~\ref{thm:char} thus reads: $\Rx$ is the union of
the medial balls and of the open half-spaces that are limits of medial
balls --- and Theorem~6 of \cite{Bia} is its faithful certificate:
such a half-space limit exists iff a supporting half-space has
non-convex contact set (Motzkin), and the level is then $0$, the drift
running over a bisector witness with no offset. The escaping class
(3) is exactly Bia\l{}o\.zyt's unreconstructed set --- the rays
leaving $\X$ with a single foot and no defective facet around them.

\paragraph{Hadamard manifolds \cite{R1,R2,R3}.} Here class (2)'s
ideal witnesses are Busemann functions at a level
(Remark~\ref{rem:offset}); the faithful certificate for their
existence is the horoconvexity defect
$\X\cap\partial\Hb{}\neq\hconv(\X\cap\partial\Hb{})$ of
\cite[Thm.~3.3--3.4]{R3}, unconditional for the $\subseteq$ direction
and achieving level $0$ in every model of the series. The hull
inclusion $\Xhat\setminus\X\subseteq\Rx$ \cite[Thm.~3.1]{R1} populates
classes (1)--(2) inside the hull; the trichotomy of
\cite[Thm.~2.1]{R3} (hull/drift/ray-escape) is
Theorem~\ref{thm:trichotomy} read through the visual boundary.

\paragraph{The sphere, compact manifolds, space forms
\cite{Sph,Tri}.} Corollary~\ref{cor:compact} contains: total
reconstruction on $\Sn$ (\cite[Thm.~4.1]{Sph}, now with the $|\X|\ge2$
hypothesis needed only for the two-point axis --- $\Rg=\Sn\setminus\X$
holds even for a single point, reconstructed from the antipode);
total reconstruction on every compact $M$
(\cite[Thm.~8.3]{Sph}, \cite[Thm.~6.4]{Tri}, caveats deleted); and,
new, total $\Rg$ on every spherical space form $\Sn/\Gamma$ for
\emph{every} closed $\X$ --- the sharp statement left open in
\cite[Rem.~7.3]{Sph}, where only the $\delta$-dense case was settled.
For the two-point axis the compact theory is unchanged:
Corollary~\ref{cor:spread} under spread, and cut-shielding as the
exact obstruction in general, with the prolate-ellipsoid transition at
eccentricity $c^{\ast}=2.6\pm0.1$ \cite[Rem.~6.6]{Tri} --- shielded
points are precisely class-(1) points of $\Rg$ whose capturing centres
lie in $\Mg\setminus\Mx$.

\paragraph{Non-compact positive curvature.} On a paraboloid of
revolution ($\sect>0$, one flattening end) with $\X$ a small circle
about the vertex, the vertex is a finite medial centre reconstructing
the inner disk (class (1)); every point outside $\X$ on the flaring
side ascends its meridian forever with a single foot and no medial
limit function reaches it: class (3), $\Rg\subsetneq M\setminus\X$.
With $\X$ two points, the bisector is unbounded and its escaping
medial centres furnish genuinely \emph{ideal} medial limit functions
on a $\sect>0$ manifold --- consistent with the trichotomy of
\cite{Tri}, since their sublevel sets are unbounded \emph{non-convex}
free regions (Busemann functions are concave above zero; it is the
maximal free \emph{convex} regions that are bounded there, not the
carriers). Partial-versus-total is decided by compactness, detection
flavour by the Hessian sign, exactly as the $2\times2$ prescribes ---
now with both rows unconditional.

\paragraph{The updated table.} Rows: compactness (axis A). Columns:
Busemann--Hessian sign (axis B). Entries: the two-geodesic theory,
with the two-point refinement in brackets.

\begin{center}
\begin{tabular}{l|ccc}
\toprule
 & $\sect\le0$ & $\sect\equiv0$ & $\sect>0$\\
\midrule
non-compact & partial; $\hconv$ defect & partial; $\conv$ defect
 & partial (paraboloid)\\
 & at level $c(\xi)$ \cite{R3} & (Bia\l{}o\.zyt \cite{Bia})
 & finite $+$ ideal centres\\
\midrule
compact & \emph{total} $\Rg$ & \emph{total} $\Rg$
 & \emph{total} $\Rg$\\
 & [$\Mx$: shielding] & [$\Mx$: shielding] & [$\Mx$: shielding;
 $=\Sn$ ok]\\
\bottomrule
\end{tabular}
\end{center}

\noindent The compact row is Corollary~\ref{cor:compact} ---
unconditional, all curvatures, $|\X|\ge1$. The non-compact row is
Theorem~\ref{thm:char}, with the faithful certificates known on the
flat and Hadamard columns and open elsewhere
(\S\ref{sec:open}).

\section{What remains}\label{sec:open}

The theory is now unconditional in its dynamical form on every
complete Riemannian manifold. What remains open is exactly the
\emph{faithful} layer --- certificates decidable from $\X$ and the
geometry of $M$, with $\Mg$ absent --- and it lives entirely on the
non-compact (escape) side.

\begin{enumerate}[label=(\arabic*)]
\item\emph{Certificates in mixed curvature.} For which ideal points
(horofunctions) $h$ of a general non-compact $M$ is $h+c$ a medial
limit function, read from the contact structure of $\X$ on
$\{h=0\}$? Known: $\conv$/Motzkin in $\R^{n}$ \cite{Bia,Mot};
$\hconv$ on Hadamard manifolds, with the Carnot and Alexandrov
frontiers of \cite[\S9--10]{R3}; vacuous on compact $M$. Unknown
already for a complete surface with both signs of curvature.
\item\emph{The level function.} On an inhomogeneous Hadamard
manifold, characterise
$c(\xi)=\inf\{c:\ b_{\xi}+c\ \text{medial limit}\}$ from $\X$ alone
(Remark~\ref{rem:offset}); $c(\xi)=0$ under pinching, in the flat
case, and over any bisector witness, and no example with
$c(\xi)>0$ is known --- deciding whether the deep-part phenomenon of
\cite[Thm.~3.4]{R3} actually occurs is the sharpest open question on
the Hadamard side.
\item\emph{The two-point axis in variable curvature.} Quantify
cut-shielding: the prolate threshold $c^{\ast}=2.6\pm0.1$ awaits an
analytic determination, and more generally the class of compact $M$
with $\Rx=M\setminus\X$ for all sparse $\X$ (constant curvature is
sufficient; bounded eccentricity conjecturally) \cite{Sph,Tri}.
\item\emph{Chebyshev problems.} The detection questions the series
isolated on its frontiers: Motzkin in Carnot groups
($\mathbb{CH}^{n}$ horospheres) and in Alexandrov spaces of
indefinite curvature (pinched horospheres) \cite[\S10]{R3}; and the
domain (Lieutier) problem on $\Sn$, where the boundary's positively
curved intrinsic geometry becomes the detection object
\cite[\S7]{Tri}.
\end{enumerate}

\paragraph{Conclusion.} Bia\l{}o\.zyt's Theorem~6 generalises to
every complete Riemannian manifold, in two layers. The
\emph{unconditional} layer is finished here: reconstruction by the
two-geodesic axis is exactly the union of the sublevel sets of the
medial limit functions (Theorem~\ref{thm:char}); every point is
captured, swallowed, or escaping (Theorem~\ref{thm:trichotomy});
capture is total on compact manifolds
(Corollary~\ref{cor:compact}), with no hypotheses whatsoever beyond
completeness and closedness. The \emph{faithful} layer --- replacing
``is a medial limit'' by a certificate on contact sets --- is exact
on the flat knife-edge ($\conv$, \cite{Bia}), on every Hadamard
manifold ($\hconv$, \cite{R3}), and trivial on compact manifolds;
its extension into mixed curvature, together with the Chebyshev-set
problems it generates, is the genuine residue. The gaps flagged in
the series --- the containment step past cut points and the radial
offset --- are closed, the first by stopping the ascent at its first
approach to the axis instead of flowing through it, the second by
normalising the limit objects by their own reach.

\begin{thebibliography}{9}

\bibitem{Bia} A.~Bia\l{}o\.zyt, \emph{Sets reconstructible with medial
axis}, preprint, AGH University of Krak\'ow, 2024.

\bibitem{R1} \emph{Sets reconstructible with medial axis: a
generalisation to Riemannian manifolds}, companion note, June 2026
(file \texttt{medial\_axis\_riemannian\_1.pdf}).

\bibitem{R2} \emph{Medial-axis reconstruction in hyperbolic space,
with a structural framework for Hadamard manifolds}, companion note,
June 2026 (file \texttt{medial\_axis\_riemannian\_2.pdf}).

\bibitem{R3} \emph{Medial-axis reconstruction in Hadamard manifolds: a
faithful, curvature-free characterization}, companion note, June 2026
(file \texttt{medial\_axis\_riemannian\_3.pdf}).

\bibitem{Sph} \emph{Medial-axis reconstruction in positive curvature:
the focal collapse and total reconstruction}, companion note, June
2026 (file \texttt{medial\_axis\_spherical.pdf}).

\bibitem{Tri} \emph{Two axes of medial-axis reconstruction:
compactness, and the Busemann--Hessian detection trichotomy},
companion note, June 2026 (file
\texttt{medial\_axis\_trichotomy.pdf}).

\bibitem{AC} J.-P.~Aubin and A.~Cellina, \emph{Differential
Inclusions}, Grundlehren 264, Springer, 1984. (Existence and extension
for upper semicontinuous convex compact right-hand sides, \S2.1.)

\bibitem{ACa} P.~Albano and P.~Cannarsa, \emph{Structural properties
of singularities of semiconcave functions}, Ann.\ Scuola Norm.\ Sup.\
Pisa \textbf{28} (1999), 719--740.

\bibitem{BH} M.~Bridson and A.~Haefliger, \emph{Metric Spaces of
Non-Positive Curvature}, Grundlehren 319, Springer, 1999.
(Horofunctions and Busemann points, II.8.)

\bibitem{Gro} M.~Gromov, \emph{Hyperbolic manifolds, groups and
actions}, in: Riemann Surfaces and Related Topics, Ann.\ of Math.\
Stud.\ 97, Princeton, 1981, 183--213. (The horofunction
compactification.)

\bibitem{Lie} A.~Lieutier, \emph{Any open bounded subset of $\R^{n}$
has the same homotopy type as its medial axis}, Computer-Aided Design
\textbf{36} (2004), 1029--1046.

\bibitem{MM} C.~Mantegazza and A.~C.~Mennucci, \emph{Hamilton--Jacobi
equations and distance functions on Riemannian manifolds}, Appl.\
Math.\ Optim.\ \textbf{47} (2003), 1--25. (Semiconcavity and first
variation of the distance from a closed set.)

\bibitem{Mot} T.~Motzkin, \emph{Sur quelques propri\'et\'es
caract\'eristiques des ensembles born\'es non convexes}, Atti Accad.\
Naz.\ Lincei \textbf{21} (1935), 773--779.

\end{thebibliography}

\end{document}
